\scanpage{22}
Cette section mérite le nom de \uline{invariant de Hasse} de $A$. L'ens des zéros de cette section est exactement l'ens des $s$ tels que $A_s$ soit \uline{non ordinaire}.

Quelles sont les relations entre $\omega_A$ et $\omega_{A^*}$ ? On a, grâce à la suite exacte
\diagram{\begin{tikzcd}[ampersand replacement=\&,column sep=small]
\unsure{0}\arrow[r,dash] \& \underline t_A^\vee\arrow[r] \& \operatorname{DR}^1(A/S)\arrow[r,dash] \& \underline t_{A^*}\arrow[r,dash] \& 0
\end{tikzcd}}
la relation
\diagram{\begin{tikzcd}[ampersand replacement=\&,column sep=large,row sep=large]
\operatorname{\text{dét}}\underline t_A^\vee\otimes\operatorname{\text{dét}}\underline t_{A^*}\arrow[r,phantom,"\simeq" description] \& \operatorname{\text{dét}}\operatorname{DR}^1(A/S)\arrow[d,phantom,"\rotatebox{90}{$\simeq$}" description] \&\\
\& \operatorname{DR}^{2g}(A/S)\arrow[r,phantom,"\simeq" description] \& \underline{\mathrm O}_S
\end{tikzcd}}
donc on aura
\[
\omega_A\cdot\omega_{A^*}^\vee\simeq\underline{\mathrm O}_S
\]
i.e. on a un isom can.
\[
\omega_A\simeq\omega_{A^*}
\]
Je dis que \uline{cet isom est compatible avec les $p^{-1}$-structures}, \uline{donc transforme $\sigma_A$ en $\sigma_{A^*}$, et dans le cas $A$ ordinaire}, induit un isom
\[
\operatorname{\text{dét}}({}_pA)_{\text{ét}}\simeq\operatorname{\text{dét}}({}_pA^*)_{\text{ét}}.
\]
\mspage{21}
